\section*{ADR-004: Extension Execution Models}
\addcontentsline{toc}{section}{ADR-004: Extension Execution Models}

\begin{description}[leftmargin=3cm,labelwidth=2.5cm]
\item[\textbf{Status}] Accepted
\item[\textbf{Date}] April 2024
\item[\textbf{Authors}] Symphony Extension Team
\end{description}

\subsection*{Summary}

We established distinct execution models for different extension types: in-process execution for infrastructure extensions (The Pit) and out-of-process execution with IPC communication for user-facing extensions (The Grand Stage). This decision defines how extensions interact with Symphony's core and each other, balancing performance requirements with safety and isolation needs.

\subsection*{Context}

Symphony's extension system must support diverse requirements across different extension categories:

\begin{itemize}
\item Infrastructure extensions requiring microsecond-level performance for AI orchestration
\item Community extensions needing safety and isolation to prevent system crashes
\item AI model integration with varying resource requirements and external dependencies
\item UI extensions with different security and performance characteristics
\item Developer tools requiring different levels of system access and privileges
\end{itemize}

The execution model choice significantly impacts system performance, reliability, security, and developer experience. Traditional approaches either sacrifice performance for safety or safety for performance, but Symphony's AI-first architecture requires both characteristics simultaneously.

\subsection*{Decision}

We implemented two distinct execution models optimized for different extension categories:

\subsubsection*{In-Process Execution (The Pit)}
Reserved for five critical infrastructure extensions running within the Conductor process:
\begin{itemize}
\item Direct memory access and function calls for maximum performance
\item Shared address space with the Conductor for zero-copy operations
\item Rust-only implementation for memory safety guarantees
\item Microsecond-level response times for critical path operations
\item Limited to trusted, thoroughly tested infrastructure components
\end{itemize}

\subsubsection*{Out-of-Process Execution (The Grand Stage)}
Used for all community-developed extensions running in isolated processes:
\begin{itemize}
\item Complete process isolation for crash protection and security
\item IPC-based communication through the Orchestra Kit
\item Support for multiple programming languages (Rust, Python, JavaScript, WASM)
\item Sandboxed execution with granular permission controls
\item Independent memory management and resource allocation
\end{itemize}

\subsection*{Rationale}

\subsubsection*{Performance Optimization}
In-process execution for The Pit enables microsecond-level response times essential for AI orchestration. Benchmarking showed 100-1000× performance improvement over IPC alternatives for infrastructure operations, enabling real-time AI model coordination.

\subsubsection*{Safety Through Isolation}
Out-of-process execution for The Grand Stage provides complete crash isolation. Extension failures, memory leaks, infinite loops, or security vulnerabilities cannot affect the core system or other extensions, ensuring system stability.

\subsubsection*{Security Model}
The dual execution model enables different security approaches: trusted infrastructure extensions with full system access, and sandboxed community extensions with restricted permissions and resource limits.

\subsubsection*{Development Flexibility}
Different execution models support different development approaches: high-performance Rust development for infrastructure, and flexible multi-language development for community extensions.

\subsubsection*{Resource Management}
In-process extensions share resources efficiently, while out-of-process extensions can be independently managed, scaled, and resource-limited without affecting other components.

\subsubsection*{Trade-offs Accepted}
\begin{itemize}
\item Increased architectural complexity requiring different development approaches
\item IPC overhead for Grand Stage extensions (0.1-0.5ms latency)
\item Memory overhead from multiple processes (50-100MB per extension)
\item Limited number of Pit extensions to maintain performance guarantees
\item Different debugging and profiling approaches for each execution model
\end{itemize}

\subsection*{Alternatives Considered}

\subsubsection*{Uniform In-Process Execution}
\textbf{Pros}: Maximum performance for all extensions, simple communication model, unified development approach, minimal memory overhead
\textbf{Cons}: Extension crashes affect entire system, security vulnerabilities propagate, difficult debugging, limited language support
\textbf{Rejected because}: Safety and reliability requirements for community extensions outweigh performance benefits

\subsubsection*{Uniform Out-of-Process Execution}
\textbf{Pros}: Complete isolation for all extensions, excellent security, independent scaling, simplified debugging
\textbf{Cons}: IPC overhead for all operations, complex state management, performance bottlenecks for infrastructure
\textbf{Rejected because}: Performance requirements for AI orchestration cannot be met with IPC overhead

\subsubsection*{Hybrid Threading Model}
\textbf{Pros}: Shared memory with thread isolation, better performance than IPC, simpler than dual processes
\textbf{Cons}: Thread crashes can affect process, complex synchronization, limited language support, security concerns
\textbf{Rejected because}: Thread-level isolation insufficient for untrusted community code

\subsubsection*{Container-Based Isolation}
\textbf{Pros}: Strong isolation, resource limits, security boundaries, language flexibility
\textbf{Cons}: High overhead, complex orchestration, startup latency, platform dependencies
\textbf{Rejected because}: Container overhead incompatible with microsecond-level performance requirements

\subsection*{Consequences}

\subsubsection*{Positive}
\begin{itemize}
\item Achieved target performance: 10-50 microseconds for Pit operations, 0.1-0.5ms for Grand Stage
\item Zero system crashes from extension failures over 8 months of operation
\item Successful community adoption with 300+ extensions developed for The Grand Stage
\item Clear architectural boundaries simplify development and maintenance
\item Flexible security model supports both trusted and untrusted code
\item Independent scaling and resource management for different extension types
\item Multi-language support enables diverse community contributions
\end{itemize}

\subsubsection*{Negative}
\begin{itemize}
\item Increased development complexity requiring expertise in both execution models
\item IPC overhead limits some real-time use cases for Grand Stage extensions
\item Memory overhead from multiple processes (average 75MB per extension)
\item Complex debugging scenarios spanning both execution environments
\item Limited promotion path from Grand Stage to Pit extensions
\item Different tooling and development workflows for each execution model
\end{itemize}

\subsection*{Success Criteria}

\begin{itemize}
\item Infrastructure operations complete within 100 microseconds (Achieved: 10-50 microseconds)
\item Grand Stage IPC latency under 1ms (Achieved: 0.1-0.5ms)
\item Zero system crashes from extension failures (Achieved: 0 crashes over 8 months)
\item Support 100+ concurrent Grand Stage extensions (Achieved: 300+ extensions)
\item Multi-language extension development support (Achieved: Rust, Python, JS, WASM)
\item Community adoption of extension development (Achieved: Active ecosystem)
\end{itemize}

\subsection*{Risks and Mitigations}

\begin{center}
\begin{tabular}{@{}p{4cm}p{2cm}p{6cm}@{}}
\toprule
\textbf{Risk} & \textbf{Impact} & \textbf{Mitigation} \\
\midrule
Architectural complexity & High & Comprehensive documentation, developer tooling, and clear guidelines \\
IPC performance bottlenecks & Medium & Optimized IPC implementation, caching, and batching strategies \\
Memory consumption & Medium & Process pooling, resource monitoring, and automatic cleanup \\
Development model confusion & Medium & Clear separation of concerns and extensive examples \\
Pit extension limitations & Low & Careful selection criteria and well-defined upgrade paths \\
Cross-process debugging & Medium & Enhanced debugging tools and logging infrastructure \\
\bottomrule
\end{tabular}
\end{center}

\subsection*{Implementation Details}

\subsubsection*{The Pit Implementation}
\begin{itemize}
\item Rust-only extensions compiled as dynamic libraries
\item Direct function calls through Foreign Function Interface (FFI)
\item Shared memory pools for zero-copy data transfer
\item Compile-time safety guarantees through Rust's type system
\item Performance monitoring and automatic failover mechanisms
\end{itemize}

\subsubsection*{The Grand Stage Implementation}
\begin{itemize}
\item Process spawning and lifecycle management through Orchestra Kit
\item JSON-RPC over multiple transport mechanisms (Unix sockets, named pipes, shared memory)
\item Sandboxed execution with configurable permission sets
\item Resource monitoring and automatic termination for misbehaving extensions
\item Hot-swappable extension updates without system restart
\end{itemize}

The dual execution model successfully enables Symphony's unique combination of high-performance AI orchestration and safe community extensibility, establishing a new paradigm for extensible AI-first development environments.